\subsection{A two-step mechanism expressed through chemical objects}
\label{sec:ni1-chemical-objects}

The Guided NI1 run computes methylamine addition to formaldehyde and carbinolamine dehydration to an imine--water complex (Figure~\ref{fig:ni1-chemical-objects}). Its Delivery DAG connects the optimized reactant complex, carbinolamine, and product complex through two first-order saddles. Calculations used B3LYP-D3(BJ)/def2-SVP in the gas phase, with charge zero and singlet multiplicity.

The addition saddle has one imaginary frequency at $-1481.06$ cm$^{-1}$ and an electronic barrier of 28.12 kcal/mol relative to the reactant complex. Its two short IRC branches lead, after endpoint optimization and frequency analysis, to the reactant complex and carbinolamine, both with zero imaginary frequencies. The dehydration saddle has one imaginary frequency at $-1796.84$ cm$^{-1}$ and an electronic barrier of 57.21 kcal/mol relative to its IRC-side carbinolamine minimum. Its two branches similarly yield carbinolamine and the imine--water complex. Each IRC followed 24 points in each direction before reaching the point limit; the endpoint calculations establish the resulting minima.

Interconversion between the carbinolamine conformers remains uncomputed. Supplementary Information Section~S4 provides the conformer references, absolute energies, native calculation records, and figure inputs.

